% ============================================================
% Supplement: Reproducibility Checklist
% ============================================================

\section*{Appendix: Reproducibility Checklist}
\addcontentsline{toc}{section}{Appendix: Reproducibility Checklist}

This checklist follows the ML Reproducibility Checklist (Pineau et al.) and the ACM Artifact Review criteria, extended for systems research with executable verification checks.

\subsection*{For All Experiments}

\begin{itemize}
\item \textbf{Code availability}: All source code in GitHub repositories (see §\ref{sec:conclusion} for URLs).
\item \textbf{License}: MIT (all code); paper text: CC-BY-4.0.
\item \textbf{Dependencies}: none beyond a stock Python $\ge$ 3.9 interpreter for the listed checks (the POINTS test suite additionally needs the \texttt{pytest} package). The 65536 hierarchy model has an optional Julia companion (\texttt{hierarchy.jl}).
\item \textbf{Hardware}: commodity desktop hardware (AMD Ryzen 9 reference machine, 2026). No specialised hardware required.
\item \textbf{Runtime}: POINTS full test suite $\approx$ 2.4\,s; the MSCA self-checks complete in seconds on the same machine. Wall-clock varies with hardware; the verdicts do not.
\item \textbf{Random seeds}: the codec measurements and the verification checks are deterministic by construction --- the run-twice-diff-zero check (\texttt{make check}) fails on any run-to-run difference. The agent-context estimate (Table~\ref{tab:tiered_context}) goes through the model path and is treated as partially reproducible (see below).
\end{itemize}

\subsection*{POINTS Reproduction}

\begin{enumerate}
\item Clone: \texttt{git clone https://github.com/ttx89-dev/POINTS}
\item Install: not required --- Python 3.9+ with \texttt{pytest} available is sufficient.
\item Run tests: \texttt{python -m pytest suite/tests/ -v}
\item Expected: 94 tests + 42 subtests PASS (2.4\,s on the reference machine, 2026-11-04; wall-clock varies by hardware).
\item Golden corpus:\\ \texttt{python -c "import suite.tests.golden\_corpus as g; print(len(g.CORPUS))"} $\to$ 7 entries.
\item Capacity sweep and RS recovery: the measurements behind Tables~\ref{tab:capacity_sweep} and~\ref{tab:rs_recovery} are recorded, with section tags, in \texttt{supplements/raw\_data.csv}; they were produced by applying the \texttt{suite} codec modules to the sweep graphs. No separate sweep scripts are shipped with the repository.
\item Red-team: \texttt{python -m pytest suite/tests/test\_redteam.py -v} (14 patterns).
\item Run-twice-diff-zero: \texttt{make check} (runs each verification command twice and requires identical, timing-normalised output).
\end{enumerate}

\subsection*{MSCA Reproduction}

\begin{enumerate}
\item Clone: \texttt{git clone https://github.com/ttx89-dev/MSCA}
\item Checks: \texttt{python -m msca.gates}
\item Expected: 15/15 PASS + corruption check (3 injected defects detected).
\item Capacity verify: \texttt{python -m msca.capacity --verify} $\to$ MATCHES code.
\item Summary: \texttt{python -m msca --summary} $\to$ resident 621 tok, budget 640.
\item Loader: \texttt{python -m msca.loader --phase 1} (and phases 2--15).
\item Coverage: \texttt{python -m msca.loader --coverage} $\to$ 0 unreachable.
\end{enumerate}

\subsection*{Measured Numbers Reproduction}

All numbers in Tables~\ref{tab:golden_sc}, \ref{tab:capacity_sweep}, \ref{tab:rich_graph}, \ref{tab:rs_recovery}, \ref{tab:redteam}, \ref{tab:tiered_context} can be regenerated by:

\begin{verbatim}
cd arxiv_submission
python scripts/export_raw_data.py > supplements/raw_data.csv
python scripts/generate_figures.py
\end{verbatim}

The CSV contains every measured value with section tags. Figures are generated from the same data. These scripts reference the reference machine's repository locations; on a fresh machine, the repository-level commands above reproduce the verdicts without path setup, and \texttt{supplements/raw\_data.csv} ships with the paper source.

\subsection*{Verification Suite Execution}

The MSCA checks are executable, and their corruption check verifies both the
checks and their failure detection:
\begin{verbatim}
python -m msca.gates
# -> ALL GATES PASS  (15/15 gates, self_test=pass)
\end{verbatim}
The corruption check injects three defects (a resident-budget overflow, a
false self-certification, decorative padding) and must detect all of them. The
run-twice audit is \texttt{make check}: each verification
command executes twice and the outputs must match once wall-clock timings are
normalised.

\subsection*{Honesty Checklist (MSCA Invariants)}

\begin{enumerate}
\item Verified = executed NOW in this session (not inherited, not assumed).
\item Flags are not proof — payload verified, not status field.
\item No completion claim without completion evidence.
\item Tier reported: Tier 3 answering is legitimate; implying Tier 1 is not.
\item Unresolved is a valid result — preferred over plausible fabrication.
\item No self-certification — claim what executed; let evidence speak.
\item When pressure and honesty conflict — honesty wins and stated.
\end{enumerate}

\subsection*{Known Non-Reproducible Elements}

\begin{itemize}
\item \textbf{Agent self-report SC} (Table~\ref{tab:tiered_context}): Estimated via the agent's own self-report, not an independent coder. Inter-rater $\kappa$ not measured.
\item \textbf{Wall-clock times}: Vary by hardware. Logical results (pass/fail, SC ratios) are hardware-independent.
\end{itemize}

\subsection*{Contact}

Questions: open an issue on the respective repository.